\documentclass{article}
\usepackage{amsmath}
\usepackage{amssymb}
\usepackage{amsthm}
\usepackage{physics}


\newtheorem*{conjecture}{Conjecture}
\newtheorem*{theorem}{Theorem}
\newtheorem*{lemma}{Lemma}


\makeatletter
\newcommand*{\rom}[1]{\expandafter\@slowromancap\romannumeral #1@}
\makeatother


\begin{document}

\section*{Exercise 2.1.1}

\subsection*{a.}

Let \((a_{1}, b_{1})\), \((a_{2}, b_{2})\), \((a_{3}, b_{3})\) form a triangle. Assume WLOG that \((a_{1}, b_{1})\) is the origin.
The are of the parallelogram defined by \((a_{2} - a_{1})\vec{i}\) + \((b_{2} - b_{1})\vec{j}\) and
\((a_{3} - a_{1})\vec{i}\) + \((b_{3} - b_{1})\vec{j}\) is:

\begin{equation*}
    \Big|\Big((a_{2} - a_{1})\vec{\imath} + (b_{2} - b_{1})\vec{\jmath} + 0\vec{k}\Big) \cross \Big((a_{3} - a_{1})\vec{\imath} + (b_{3} - b_{1})\vec{\jmath} + 0\vec{k}\Big)\Big| =
\end{equation*}
\begin{equation*}
    \abs{
\begin{vmatrix} 
    \vec{\imath} & \vec{\jmath} & \vec{k} \\
    (a_{2} - a_{1}) & (b_{2} - b_{1}) & 0 \\
    (a_{3} - a_{1}) & (b_{3} - b_{1}) & 0
\end{vmatrix}
    }
    =
\end{equation*}
\begin{multline*}
        = \Bigg|\Big( (b_{2} - b_{1}) \cdot 0 - 0 \cdot (b_{3} - b_{1}) \Big)\vec{\imath} - \Big( (a_{2} - a_{1}) \cdot 0 - 0 \cdot (a_{3} - a_{1}) \Big)\vec{\jmath} \\
        + \Big( (a_{2} - a_{1}) \cdot (b_{3} - b_{1}) - (b_{2} - b_{1}) \cdot (a_{3} - a_{1}) \Big)\vec{k}\Bigg| =
\end{multline*}
\begin{equation*}
    = \abs{\Big( (a_{2} - a_{1}) \cdot (b_{3} - b_{1}) - (b_{2} - b_{1}) \cdot (a_{3} - a_{1}) \Big)\vec{k}}
\end{equation*}
\begin{equation*}
    = \Big| (a_{2} - a_{1}) \cdot (b_{3} - b_{1}) - (b_{2} - b_{1}) \cdot (a_{3} - a_{1})\Big| =
\end{equation*}


Since the parallelogram is built from the triangle with the defined vertices, it can equivalently be seen built from two
triangles with the same area. Hence the are of the triangle \((a_{1}, b_{1})\), \((a_{2}, b_{2})\), \((a_{3}, b_{3})\) is:

\begin{equation*}
    \frac{1}{2} \Big| (a_{2} - a_{1}) \cdot (b_{3} - b_{1}) - (b_{2} - b_{1}) \cdot (a_{3} - a_{1})\Big|
\end{equation*}


\subsection*{b.}

Let vertices \((a, a^{2})\), \((a + \delta, (a + \delta)^{2})\), \((a + 2\delta, (a + 2\delta)^{2})\) form a
a triangle. Then using the formula from point (a):

\begin{equation*}
    A = \frac{1}{2}\Big|(a + \delta - a)\Big((a + 2\delta)^{2} - a^{2}\Big) - (\Big(a + \delta\Big)^{2} - a^{2})(a + 2\delta - a)\Big| =
\end{equation*}
\begin{equation*}
    = \frac{1}{2}\Big|\delta(a^{2} + 4a\delta + 4\delta^{2} - a^{2}) - (a^{2} + 2a\delta + \delta^{2} - a^{2})(2\delta)\Big|
\end{equation*}
\begin{equation*}
    = \frac{1}{2}\Big|4a\delta^{2} + 4\delta^{3} - (4a\delta^{2} + 2\delta^{3})\Big|
\end{equation*}
\begin{equation*}
    = \frac{1}{2}\Big|4\delta^{3} - 2\delta^{3}\Big| = \Big|\delta\Big|^{3}
\end{equation*}

\subsubsection*{Errata from Bressoud}

Let vertices \((a, 1 - a^{2})\), \((a + \delta, 1 - (a + \delta)^{2})\), \((a + 2\delta, 1 - (a + 2\delta)^{2})\) form a
a triangle. Then using the formula from point (a):

\begin{equation*}
    A = \frac{1}{2}\Big|(a + \delta - a)\Big(1 - (a + 2\delta)^{2} - 1 + a^{2}\Big) - \Big((1 - (a + \delta)^{2} - 1 + a^{2})(a + 2\delta - a)\Big)\Big| =
\end{equation*}
\begin{equation*}
    = \frac{1}{2}\Big|\delta(-a^{2} - 4\delta a - 4\delta^{2} + a^{2}) - \Big((-a^{2} - 2a\delta - \delta^{2} + a^{2})2\delta\Big)\Big|
\end{equation*}
\begin{equation*}
    = \frac{1}{2}\Big|\delta(-4\delta a - 4\delta^{2}) - \Big((- 2a\delta - \delta^{2})2\delta\Big)\Big|
\end{equation*}
\begin{equation*}
    = \frac{1}{2}\Big|-4\delta^{2}a - 4\delta^{3} - (-4a\delta^{2} - 2\delta^{3})\Big| =
\end{equation*}
\begin{equation*}
    = \frac{1}{2}\Big|-4\delta^{2}a - 4\delta^{3} + 4a\delta^{2} + 2\delta^{3}\Big| = 
\end{equation*}
\begin{equation*}
    = \frac{1}{2}\Big| -4\delta^{3} + 2\delta^{3} \Big| = \frac{1}{2}\Big| -2\delta^{3} \Big| = 
\end{equation*}
\begin{equation*}
    = \frac{1}{2}|-2||\delta^{3}| = |\delta^{3}| = |\delta|^{3}
\end{equation*}


\subsection*{c.}




\end{document}